% !TeX root = ../../Book2.tex
\section{泛函与变分}
\label{b2:sec:1401}

固定两端的曲线可以改变形状，而端点条件保持不变。若用函数 \(u\) 表示曲线，一次允许的微小改变便写成 \(u+th\)，其中 \(h\) 在端点为零。长度、面积和能量都是作用于整条曲线的函数；考察它们随 \(t\) 的变化，就把极值问题带回了一元微积分。不过，在函数空间中，沿每个方向分别求导与对所有小扰动作统一估计并不是同一件事。本节先辨明这一差别，再说明积分泛函的变分怎样计算。

\subsection{Gâteaux 导数}
\label{b2:subsec:140101}

设 \(X\) 是实 Banach 空间，\(U\subset X\) 是开集，\(J:U\to\mathbb R\)。给定 \(u\in U\) 和 \(h\in X\)，若极限存在，记
\[
 \delta J(u;h)=\lim_{t\to0}\frac{J(u+th)-J(u)}{t}.
\]
这里的参数 \(t\) 取实数；当未知函数取复值而泛函取实值时，也应把函数空间看作实线性空间。比如 \(J(u)=\lVert u\rVert^2\) 的变分含有内积的实部，并不是复线性泛函。

\Needspace{6\baselineskip}
\begin{definition}[Gâteaux 可微]
\label{b2:def:14-gateaux}
设 \(X\) 是实 Banach 空间，\(J:U\to\mathbb R\)，\(U\subset X\) 开、\(u\in U\)；\(\delta J(u;h)=\lim_{t\to0}(J(u+th)-J(u))/t\)。若各方向的变分都存在，且存在连续线性泛函 \(A\in X^*\)，使得 \(\delta J(u;h)=Ah\) 对所有 \(h\in X\) 成立，则称 \(J\) 在 \(u\) 处 \term[gateaux-kewei]{Gâteaux 可微}[Gâteaux differentiable]，并称 \(A\) 为其 \term[gateaux-daoshu]{Gâteaux 导数}[Gâteaux derivative]，记作 \(D_GJ(u)\)。
\end{definition}

定义中的线性性不是由方向极限自动保证的。在 \(\mathbb R^2\) 上令
\[
 F(x,y)=\begin{cases}x^3/(x^2+y^2),&(x,y)\ne(0,0),\\0,&(x,y)=(0,0).\end{cases}
\]
沿方向 \((a,b)\ne(0,0)\) 的导数是 \(a^3/(a^2+b^2)\)。这个表达式关于方向不是线性的，故 \(F\) 虽有所有方向导数，却不满足这里的 Gâteaux 可微定义。不同文献有时把“方向导数均存在”也称为 Gâteaux 可微；使用结论时，应核对是否同时要求导数为连续线性映射。

\Needspace{6\baselineskip}
\begin{proposition}[极小点的必要条件]
\label{b2:prop:14-gateaux-minimum}
设 \(X\) 是实 Banach 空间，\(V\subset X\) 是线性子空间，\(J\) 是定义在 \(u\in X\) 某开邻域内的实值泛函。若 \(u\) 是 \(J\) 在仿射空间 \(u+V\) 上的局部极小点，且 \(\delta J(u;h)\) 对每个 \(h\in V\) 都存在，则
\[\delta J(u;h)=0\qquad(h\in V).\]
\end{proposition}
\begin{proof}
固定 \(h\in V\)。当 \(t\) 的绝对值充分小时，\(J(u+th)-J(u)\ge0\)。除以正的 \(t\) 再令 \(t\downarrow0\)，得变分非负；除以负的 \(t\) 再令 \(t\uparrow0\)，得变分非正。两个单侧极限相同，故变分为零。
\end{proof}

此处真正起作用的是正负两个方向都可行。若要求 \(u\ge0\)，在接触零值的区域只能作某些单侧扰动，必要条件一般变成不等式。类似地，固定端点要求扰动在端点消失；固定积分则要求扰动的积分为零。求第一变分之前，先写清允许的扰动，比形式地求导更为重要。

\subsection{Fréchet 导数}
\label{b2:subsec:140102}

第一章在有限维空间中讨论过 Fréchet 导数；同一定义适用于 Banach 空间。设 \(X\) 为实 Banach 空间，\(U\subset X\) 开，\(J:U\to\mathbb R\)、\(u\in U\)。泛函 \(J\) 在 \(u\) 处 Fréchet 可微，意味着存在 \(A\in X^*\)，使
\[
 J(u+h)=J(u)+Ah+r(h),\qquad
 \frac{|r(h)|}{\lVert h\rVert_X}\longrightarrow0\quad(h\to0).
\]
这里的余项估计对所有小方向统一成立。因此 Fréchet 可微必然蕴含 Gâteaux 可微，且两种导数相同；反过来则需要额外条件。

\Needspace{6\baselineskip}
\begin{proposition}[导数连续时的统一逼近]
\label{b2:prop:14-continuous-gateaux}
设 \(X\) 是实 Banach 空间，实值泛函 \(J\) 在 \(u\) 的一个开邻域内 Gâteaux 可微。若 \(v\mapsto D_GJ(v)\) 在 \(u\) 处关于 \(X^*\) 的算子范数连续，则 \(J\) 在 \(u\) 处 Fréchet 可微，且 \(DJ(u)=D_GJ(u)\)。
\end{proposition}
\begin{proof}
给定 \(\varepsilon>0\)，选取 \(r>0\)，使 \(\lVert v-u\rVert_X<r\) 时有 \(\lVert D_GJ(v)-D_GJ(u)\rVert_{X^*}<\varepsilon\)。对 \(\lVert h\rVert_X<r\)，在 \([0,1]\) 上考察
\[g(t)=J(u+th)-J(u)-tD_GJ(u)h.\]
它在区间上连续、内部可微，且 \(|g'(t)|\le\varepsilon\lVert h\rVert_X\)。一元中值定理给出 \(|g(1)|\le\varepsilon\lVert h\rVert_X\)，正是所需的统一余项估计。
\end{proof}

即便 Gâteaux 导数在一点是零映射，原函数也未必在该点连续。\subsecref{b2:subsec:010101}已经计算过的反例是
\[
 J(x,y)=\begin{cases}\displaystyle\frac{x^3y}{x^6+y^2},&(x,y)\ne(0,0),\\0,&(x,y)=(0,0).\end{cases}
\]
它在原点的所有方向导数均为零，即 \(D_GJ(0)=0\)，但沿 \(y=x^3\)、\(x\ne0\) 有 \(J(x,x^3)=1/2\)，故原点处不连续，也不可能 Fréchet 可微。这里需要保留的新区别是：每个固定方向的极限存在，并不提供对小扰动统一的余项估计；前一命题增加的导数算子范数连续性，才足以把这些逐方向极限接成同一个局部线性逼近。

\subsection{第一变分}
\label{b2:subsec:140103}

对实 Banach 空间 \(X\) 的开集 \(U\) 上的 Fréchet 可微泛函 \(J:U\to\mathbb R\)，\(u\in U\)、\(h\in X\) 时的 \(DJ(u)h\) 即泛函对扰动的一阶线性响应，称为 \term[diyi-bianfen]{第一变分}[first variation]。对积分泛函，常见计算是“把导数移入积分号”。这一计算需要一致的支配或余项估计；下面先在 \(C^1\) 空间内证明一个便于直接使用的版本。

\Needspace{6\baselineskip}
\begin{proposition}[光滑积分泛函的第一变分]
\label{b2:prop:14-smooth-integral-variation}
设整数 \(d\geq1\)，\(\Omega\subset\mathbb R^d\) 是有界开集。实函数 \(F(x,z,p)\) 以及对 \((z,p)\in\mathbb R\times\mathbb R^d\) 的一阶偏导数在 \(\overline\Omega\times\mathbb R\times\mathbb R^d\) 上连续。赋予实函数空间 \(C^1(\overline\Omega;\mathbb R)\) 以函数及一阶导数的一致范数。则
\[J(u)=\int_\Omega F(x,u(x),\nabla u(x))\,\mathrm dx\]
为 Fréchet 可微泛函，且
\[
 DJ(u)h=\int_\Omega\bigl(F_z(x,u,\nabla u)h+
             F_p(x,u,\nabla u)\cdot\nabla h\bigr)\,\mathrm dx.
\]
其导数关于 \(u\) 在算子范数下连续。
\end{proposition}
\begin{proof}
固定 \(u\)。当 \(\lVert h\rVert_{C^1}\le1\) 且 \(0\le t\le1\) 时，\((x,u(x)+th(x),\nabla u(x)+t\nabla h(x))\) 落在一个固定紧集内。在这个紧集上，\((z,p)\) 的一阶偏导数一致连续；以 \(\omega(s)\) 表示相应连续模，其中 \(\omega(s)\to0\)。沿 \((h,\nabla h)\) 应用一元微积分基本定理，得
\[
 \begin{split}
 &F(x,u+h,\nabla u+\nabla h)-F(x,u,\nabla u)\\
 &\quad-F_z(x,u,\nabla u)h-F_p(x,u,\nabla u)\cdot\nabla h\\
 &=\int_0^1\bigl[D_{(z,p)}F(x,u+th,\nabla u+t\nabla h)
              -D_{(z,p)}F(x,u,\nabla u)\bigr]\cdot(h,\nabla h)\,\mathrm dt.
 \end{split}
\]
积分后，余项的绝对值不超过
\(C_d|\Omega|\cdot\omega(C_d\lVert h\rVert_{C^1})\cdot\lVert h\rVert_{C^1}\)。所写的一阶项是有界线性泛函，故它就是 Fréchet 导数。用同一个紧集上的连续模比较 \(DJ(u+k)\) 与 \(DJ(u)\)，再对 \(\lVert h\rVert_{C^1}\le1\) 取上确界，就得到导数的算子范数连续性。
\end{proof}

例如，在有界开集 \(\Omega\) 上取实值 \(f\in C(\overline\Omega)\)，则 \(F(x,z,p)=\tfrac12|p|^2-f(x)z\) 满足前一命题的连续性条件，在 \(C^1(\overline\Omega;\mathbb R)\) 上给出
\[DJ(u)h=\int_\Omega\nabla u\cdot\nabla h\,\mathrm dx-\int_\Omega fh\,\mathrm dx.\]
改在 \(H^1(\Omega)\) 上讨论时，上一命题不能原样使用：\(H^1\) 范数一般不控制函数的一致范数，非线性项的增长速度必须与 Sobolev 嵌入相容。特别是，写下 \(\int |u|^q\) 之前，须先确认每个允许的 \(u\) 确有这一可积性。

二次能量可以避开一致范数，直接在其自然空间中计算。设 \(\Omega\subset\mathbb R^d\) 为开集，只要求实值 \(f\in L^2(\Omega)\)，令
\[
 J_f(u)=\frac12\int_\Omega|\nabla u|^2\,\mathrm{d}x
          -\int_\Omega fu\,\mathrm{d}x,
 \qquad u\in H^1(\Omega;\mathbb R).
\]
对每个 \(u,h\in H^1(\Omega;\mathbb R)\)，候选线性响应
\[
 A_uh=\int_\Omega\nabla u\cdot\nabla h\,\mathrm{d}x
             -\int_\Omega fh\,\mathrm{d}x
\]
满足 \(|A_uh|\leq(\|\nabla u\|_2+\|f\|_2)\|h\|_{H^1}\)。直接展开平方还给出精确余项
\[
 J_f(u+h)-J_f(u)-A_uh=\frac12\|\nabla h\|_2^2
                       \leq\frac12\|h\|_{H^1}^2.
\]
除以 \(\|h\|_{H^1}\) 后趋零，故 \(DJ_f(u)=A_u\)。而 \(\|A_u-A_v\|_{(H^1)^*}\leq\|\nabla(u-v)\|_2\)，导数在算子范数下连续。这里的 \(f\) 不必连续；\(L^2\) 配对的有界性及二次余项已经给出所需的可微性。本章后面将以这个空间中的能量进行极小化。

梯度幂能量可以直接在自然的可积空间中微分。这里不能用逐点求导代替范数余项估计，因为扰动可能在很小的集合上很大。

\Needspace{6\baselineskip}
\begin{proposition}[梯度幂能量的连续微分]
\label{b2:prop:14-power-frechet}
设整数 \(d\geq1\)，\(\Omega\subset\mathbb R^d\) 为开集。若 \(1<p<\infty\)，则 \(P(u)=p^{-1}\int_\Omega|\nabla u|^p\) 在实空间 \(W^{1,p}(\Omega)\) 上连续 Fréchet 可微，且
\[DP(u)h=\int_\Omega |\nabla u|^{p-2}\nabla u\cdot\nabla h.\]
\end{proposition}
\begin{proof}
置 \(A_p(z)=|z|^{p-2}z\)，并规定 \(A_p(0)=0\)；\(p=2\) 时直接有 \(A_2(z)=z\)。先核验使积分余项一致变小的两个向量估计：
\[
 |A_p(a+b)-A_p(a)|\le
 \begin{cases}
 C_p(|a|+|b|)^{p-2}|b|,&p\ge2,\\
 C_p|b|^{p-1},&1<p<2.
 \end{cases}
\]
在非零点，\(DA_p(z)=|z|^{p-2}I+(p-2)|z|^{p-4}z\otimes z\)，其算子范数不超过 \(C_p|z|^{p-2}\)。当 \(p\ge2\) 时，沿 \(a+tb\) 积分即得第一式，经过零点也可分成两段取极限。当 \(p<2\)、\(|b|\ge |a|/2\) 时，分别估计两端的 \(|A_p|\)，得到第二式；当 \(|b|<|a|/2\) 时，整条线段离零点至少 \(|a|/2\)，导数估计给出 \(C_p|a|^{p-2}|b|\le C_p|b|^{p-1}\)。

对 \(v,w\in L^p(\Omega;\mathbb R^d)\)，Hölder 不等式于是给出
\[
 \|A_p(v+w)-A_p(v)\|_{p'}\le
 \begin{cases}
 C_p(\|v\|_p+\|w\|_p)^{p-2}\|w\|_p,&p\ge2,\\
 C_p\|w\|_p^{p-1},&1<p<2,
 \end{cases}
 \qquad p'=\frac p{p-1}.
\]
第一种情形用的指数关系是 \(1/p'=(p-2)/p+1/p\)；第二种情形直接把 \(|w|^{p-1}\) 的 \(p'\) 次方积分。另有 \(\|A_p(v)\|_{p'}=\|v\|_p^{p-1}\)，故候选导数确为有界线性泛函。

沿逐点线段积分，再用 Tonelli 与 Hölder，余项满足
\[
 |P(u+h)-P(u)-DP(u)h|
 \le \|\nabla h\|_p\sup_{0\le t\le1}
 \|A_p(\nabla u+t\nabla h)-A_p(\nabla u)\|_{p'}.
\]
上界在 \(p\ge2\) 时为 \(O((\|\nabla u\|_p+\|\nabla h\|_p)^{p-2}\|\nabla h\|_p^2)\)，在 \(p<2\) 时为 \(O(\|\nabla h\|_p^p)\)，除以 \(\|h\|_{W^{1,p}}\) 后都趋零。同一向量估计又说明 \(DP(u+h)\to DP(u)\) 于对偶算子范数，完成证明。
\end{proof}

这也辨明了第一与第二变分的差别：\(1<p<2\) 时一阶导数连续，但 \(|z|^p/p\) 的二阶导数在 \(z=0\) 附近发散；不能把二次能量的有界 Hessian 论证照搬到所有幂能量。下面的存在性只需一阶变分和凸性，不需这种额外二阶正则性。

\subsection{第二变分}
\label{b2:subsec:140104}

设 \(X\) 为实 Banach 空间，\(U\subset X\) 开，\(J:U\to\mathbb R\)，\(u\in U\)、\(h\in X\)，且 \(J\) 在 \(u\) 的邻域可微。当导数映射本身在 \(u\) 处 Fréchet 可微时，其导数 \(D^2J(u)\) 是有界双线性型。对同一方向两次求导所得的
\[\delta^2J(u;h)=D^2J(u)[h,h]\]
称为 \term[dier-bianfen]{第二变分}[second variation]。若 \(J\) 在 \(u\) 附近二阶连续 Fréchet 可微，则沿线段积分两次得到
\[
 J(u+h)-J(u)-DJ(u)h
 =\int_0^1(1-t)D^2J(u+th)[h,h]\,\mathrm dt.
\]
二阶项反映驻点附近的弯曲方向；它提供必要条件，也在具有统一正下界时提供充分条件。

\Needspace{6\baselineskip}
\begin{theorem}[第二变分判别]
\label{b2:thm:14-second-variation-test}
设 \(X\) 是实 Banach 空间，\(u\in X\)，\(V\subset X\) 是闭线性子空间，实值泛函 \(J\) 在 \(u\) 的开邻域内二阶连续 Fréchet 可微。若 \(u\) 是 \(J\) 在 \(u+V\) 上的局部极小点，则
\[DJ(u)h=0,\qquad D^2J(u)[h,h]\ge0\quad(h\in V).\]
反之，若 \(DJ(u)|_V=0\)，且存在 \(c>0\)，使
\[D^2J(u)[h,h]\ge c\lVert h\rVert_X^2\quad(h\in V),\]
则 \(u\) 是 \(J\) 在 \(u+V\) 上的严格局部极小点。
\end{theorem}
\begin{proof}
必要性由一元函数 \(t\mapsto J(u+th)\) 的一、二阶极值条件得到。证明充分性时，利用 \(D^2J\) 的算子范数连续性，取 \(r>0\)，使 \(\lVert v-u\rVert_X<r\) 时 \(\lVert D^2J(v)-D^2J(u)\rVert<c/2\)。对 \(h\in V\)、\(0<\lVert h\rVert_X<r\)，上面的积分余项公式给出
\[
 J(u+h)-J(u)\ge\int_0^1(1-t)\frac c2\lVert h\rVert_X^2\,\mathrm dt
 =\frac c4\lVert h\rVert_X^2>0.
\]
\end{proof}

仅有非负第二变分并不足以判别极小：\(t^4\) 与 \(-t^4\) 在零点的第一、第二导数都为零，却有相反的极值性质。在无限维空间里，即使对每个非零 \(h\) 都有 \(D^2J(u)[h,h]>0\)，也不能直接得到定理要求的统一常数 \(c\)。后面处理 Dirichlet 能量时，正是 Poincaré 不等式把梯度能量转成整个函数空间范数的控制，从而弥补这一差别。
